%% =====================================================================
%%  B3-T-02-01 -- what B3 contributes to "Affect Control"
%%  -------------------------------------------------------------------
%%  LAYER A: APPEND-ONLY. Written once, then left alone. Material found
%%  only here sits in the `unique` slot and is protected by rule R10.
%% =====================================================================
%% @kind: source
%% @id: B3-T-02-01
%% @book: B3
%% @topic: T-02-01
%% @locator: Law 39, pp. 325--332; Law 43
%% @status: distilled
%% @unique: yes
%% @integrated: yes
%% @updated: 2026-09-19

\dossier{B3}{T-02-01}{\bookshort{B3}}{Law 39, pp. 325--332; Law 43}

\begin{distilled}
  \item Anger and emotion are treated as strategically counterproductive for the person feeling them; the instruction is to stay calm and objective at all times.
  \item The same law turns offensive: making an opponent angry while remaining calm yourself is described as a decisive advantage, obtained by finding the thread in their vanity.
  \item Working on hearts and minds is preferred to coercion because coerced compliance generates a reaction that later turns against the user.
\end{distilled}

\begin{unique}
  \item The two-directional treatment of affect: self-composure and provoking the opponent's emotion are presented as one technique with two ends, and vanity is named as the usual thread.
\end{unique}
